% 第九卷 3.6 节：四种目标配套的可见范围（注意力掩码）对照。
% 行为查询位置 i，列为键位置 j；六个位置依次为“读书会／改到／周五／下午／两点／举行”。
% 蓝格：同一组参数内的自注意力可见；黄格：编码器—解码器中经交叉注意力读取；白格：不可见。
\begin{tikzpicture}[
  x=1cm,
  y=1cm,
  font=\figurefont\footnotesize,
  seen/.style={draw=white,line width=0.5pt,fill=FigureBlue!70},
  cross/.style={draw=white,line width=0.5pt,fill=FigureYellow!85},
  hidden/.style={draw=FigureGrid,line width=0.4pt,fill=white},
  idx/.style={font=\figurefont\scriptsize,text=FigureMuted},
  ptitle/.style={font=\figurefont\bfseries\footnotesize,text=FigureInk,anchor=south},
  note/.style={font=\figurefont\scriptsize,text=FigureMuted,anchor=north,align=center}
]
  \def\s{0.40}
  % 面板编号 \p：1 因果，2 前缀（P=2），3 双向，4 编码器—解码器（源为位置 1–3）。
  \foreach \p/\xo/\ttl/\nt in {%
      1/0.0/{(a) 因果（下一词元预测）}/{位置$i$只读$j\leq i$},%
      2/3.55/{(b) 前缀语言模型}/{前缀$P=2$内互读，\\其后因果},%
      3/7.1/{(c) 双向（掩码语言模型）}/{全部互读；\\被遮位置输入为掩码},%
      4/10.65/{(d) 编码器—解码器}/{源1--3双向；目标4--6因果，\\并经交叉注意力读源}} {
    \foreach \i in {1,...,6} {
      \foreach \j in {1,...,6} {
        \pgfmathtruncatemacro{\kind}{%
          \p==1 ? (\j<=\i ? 1 : 0) :
          (\p==2 ? ((\j<=2 || \j<=\i) ? 1 : 0) :
          (\p==3 ? 1 :
          (\i<=3 ? (\j<=3 ? 1 : 0) : (\j<=3 ? 2 : (\j<=\i ? 1 : 0)))))}
        \ifcase\kind
          \path[hidden] ({\xo+(\j-1)*\s},{-(\i-1)*\s}) rectangle ++(\s,-\s);
        \or
          \path[seen] ({\xo+(\j-1)*\s},{-(\i-1)*\s}) rectangle ++(\s,-\s);
        \or
          \path[cross] ({\xo+(\j-1)*\s},{-(\i-1)*\s}) rectangle ++(\s,-\s);
        \fi
      }
      \node[idx,anchor=east] at ({\xo-0.04},{-(\i-0.5)*\s}) {\i};
      \node[idx,anchor=south] at ({\xo+(\i-0.5)*\s},0.02) {\i};
    }
    \draw[FigureStroke,line width=0.6pt] (\xo,0) rectangle ++({6*\s},{-6*\s});
    \node[ptitle] at ({\xo+3*\s},0.36) {\ttl};
    \node[note] at ({\xo+3*\s},{-6*\s-0.12}) {\nt};
  }
  % 前缀与源的分界线。
  \draw[FigureStroke,line width=0.9pt,dashed] ({3.55+2*\s},0) -- ({3.55+2*\s},{-6*\s});
  \draw[FigureStroke,line width=0.9pt,dashed] (3.55,{-2*\s}) -- ({3.55+6*\s},{-2*\s});
  \draw[FigureStroke,line width=0.9pt] ({10.65+3*\s},0) -- ({10.65+3*\s},{-6*\s});
  \draw[FigureStroke,line width=0.9pt] (10.65,{-3*\s}) -- ({10.65+6*\s},{-3*\s});
  % 坐标说明与图例。
  \node[idx,anchor=east] at (-0.3,{-3*\s}) {查询$i$};
  \node[idx,anchor=south east] at (-0.06,0.02) {键$j$};
  \path[seen] (3.2,-3.75) rectangle ++(\s,-\s);
  \node[anchor=west,font=\figurefont\scriptsize,text=FigureInk] at ({3.2+\s+0.08},{-3.75-0.5*\s}) {自注意力可见};
  \path[cross] (6.0,-3.75) rectangle ++(\s,-\s);
  \node[anchor=west,font=\figurefont\scriptsize,text=FigureInk] at ({6.0+\s+0.08},{-3.75-0.5*\s}) {交叉注意力读取};
  \path[hidden] (8.9,-3.75) rectangle ++(\s,-\s);
  \node[anchor=west,font=\figurefont\scriptsize,text=FigureInk] at ({8.9+\s+0.08},{-3.75-0.5*\s}) {不可见};
\end{tikzpicture}
